% =====================================================================
% Capítulo 2 — Marco teórico
% =====================================================================
\chapter{Marco teórico: fundamentos de los modelos de difusión}
\label{cap:marco-teorico}

% ---------------------------------------------------------------------
\section{¿Qué es un modelo de difusión?}
\label{sec:que-es-difusion}

Un \emph{modelo de difusión} es un modelo generativo probabilístico
que aprende a producir muestras nuevas de una distribución de datos
$p(x)$ invirtiendo, paso a paso, un proceso de corrupción gradual.
La idea central es deliberadamente sencilla: si se sabe cómo destruir
la estructura de un dato añadiéndole ruido de forma progresiva hasta
convertirlo en ruido puro, entonces aprender a \emph{deshacer} cada
uno de esos pasos equivale a aprender a \emph{crear} datos a partir de
ruido. El modelo se entrena, por tanto, no para describir
explícitamente la distribución $p(x)$ —algo intratable para datos de
alta dimensión como las imágenes— sino para revertir un proceso de
difusión conocido.

Conviene fijar desde el principio los dos ingredientes que
reaparecerán a lo largo de todo el capítulo:

\begin{itemize}
    \item El \textbf{proceso hacia delante} (\textit{forward} o
    \textit{diffusion process}): una cadena fija, sin parámetros
    entrenables, que añade ruido gaussiano a un dato $x_0$ a lo largo
    de $T$ pasos hasta que $x_T$ es indistinguible de una muestra de
    $\mathcal{N}(0,\mathbf{I})$.
    \item El \textbf{proceso hacia atrás} (\textit{reverse} o
    \textit{denoising process}): la cadena que queremos aprender, en
    la que una red neuronal con parámetros $\theta$ estima cómo
    eliminar el ruido en cada paso para reconstruir progresivamente
    una muestra de datos a partir de ruido puro.
\end{itemize}

A diferencia de los GAN, donde un generador y un discriminador se
entrenan de forma adversaria, o de los VAE, que comprimen los datos
en un espacio latente de baja dimensión, los modelos de difusión
descomponen la generación en una secuencia larga de pasos de
\emph{denoising} pequeños y bien condicionados, cada uno de los cuales
es un problema de aprendizaje supervisado sencillo y estable. Esta
estabilidad de entrenamiento, junto con la calidad y diversidad de las
muestras, es lo que ha convertido a la difusión en el paradigma
dominante de la IA generativa actual.

\todo{Decidir si se incluye aquí una figura-resumen que ilustre las
dos cadenas (forward/reverse) sobre una imagen de ejemplo, o si se
reutiliza la figura del capítulo de mariposas. Añadir comparativa
breve con GAN/VAE si se quiere reforzar el posicionamiento.}

% ---------------------------------------------------------------------
\section{¿Por qué pareció interesante la difusión?}
\label{sec:por-que-difusion}

\subsection{El origen físico: termodinámica del no equilibrio}
\label{subsec:origen-fisico}

La motivación original de los modelos de difusión es explícitamente
física. El trabajo fundacional de Sohl-Dickstein et
al.~\cite{sohldickstein2015} se inspira en la \emph{termodinámica del
no equilibrio}: del mismo modo que una gota de tinta en un vaso de
agua se difunde hasta repartirse de manera homogénea —un estado de
máxima entropía, fácil de describir—, una distribución de datos
compleja puede transformarse, mediante un proceso de difusión, en una
distribución gaussiana simple. La apuesta teórica fue que, si dicho
proceso de difusión se aplica con pasos suficientemente pequeños,
cada paso de la \emph{inversión} de ese proceso admite también una
forma funcional sencilla (gaussiana), lo que lo hace aprendible por
una red neuronal. Es decir: se renuncia a modelar directamente la
complejísima $p(x)$ y, a cambio, se modela una larga secuencia de
transiciones locales, cada una de ellas tratable.

Esta analogía con la física no es meramente decorativa. La conexión
formal con las ecuaciones de difusión y, más adelante, con las
ecuaciones diferenciales estocásticas (SDE), dota al método de un
aparato matemático maduro —el de los procesos de difusión y la
ecuación de Fokker-Planck— del que se pueden tomar prestadas
herramientas de análisis y de muestreo.

\subsection{Score-matching: la otra raíz teórica}
\label{subsec:score-matching}

Existe una segunda línea de trabajo, inicialmente independiente, que
conviene distinguir con cuidado porque resulta clave para responder a
la pregunta de \emph{por qué} la difusión pareció prometedora: el
\emph{score-matching}. En esta formulación, propuesta por Song y
Ermon~\cite{song2019scorematching}, el objetivo no es modelar la
densidad $p(x)$ sino su \emph{score}, esto es, el gradiente del
logaritmo de la densidad respecto a los datos,
$\nabla_x \log p(x)$. Conocer este campo vectorial basta para generar
muestras: partiendo de ruido, se siguen iterativamente los gradientes
hacia regiones de alta densidad mediante \emph{dinámica de Langevin}.

El problema práctico es que el score está mal estimado en las zonas de
baja densidad, precisamente donde el muestreo arranca. La solución que
propusieron —perturbar los datos con ruido gaussiano a múltiples
escalas y estimar el score de cada versión ruidosa— es,
sorprendentemente, casi idéntica en la práctica al proceso de
difusión de Sohl-Dickstein. Esta convergencia se formalizó después en
el marco unificado de las ecuaciones diferenciales estocásticas de
Song et al.~\cite{song2021sde}, donde tanto los DDPM como los modelos
de score-matching se revelan como dos discretizaciones distintas de la
misma SDE subyacente.

\subsection{¿Y todo esto aplica a los modelos que usamos hoy?}
\label{subsec:aplica-hoy}

Esta es la pregunta que motiva el matiz más importante del capítulo.
La intuición física de la difusión y el marco de score-matching son
\emph{dos perspectivas teóricas distintas sobre el mismo objeto}, y
ambas se aplican —en el sentido de que ofrecen una explicación
válida— a los modelos que utilizamos en este trabajo. Sin embargo, es
importante no confundir la \emph{justificación teórica} con la
\emph{formulación operativa}:

\begin{itemize}
    \item Los modelos que empleamos aquí pertenecen a la familia
    \textbf{DDPM}~\cite{ho2020ddpm} y sus descendientes
    deterministas (DDIM~\cite{song2021ddim}). Su entrenamiento no se
    plantea explícitamente como score-matching ni como la inversión
    de una SDE, sino como un objetivo de \emph{predicción del ruido}:
    la red recibe una muestra ruidosa $x_t$ y el paso $t$, y aprende
    a predecir el ruido $\epsilon$ que se le añadió.
    \item No obstante, puede demostrarse que ese objetivo de
    predicción de ruido es \emph{equivalente}, salvo una constante de
    reescalado dependiente de $t$, a estimar el score
    $\nabla_{x_t}\log q(x_t)$ de la distribución ruidosa. Es decir:
    aunque los DDPM no se \emph{deriven} desde el score-matching, sí
    lo \emph{realizan} implícitamente.
\end{itemize}

En resumen, la respuesta a la pregunta planteada es matizada: la raíz
física-termodinámica y el score-matching no son frameworks rivales
sino caras de una misma teoría, y ambos describen correctamente lo que
hacen los modelos modernos; pero la maquinaria concreta con la que
entrenamos e inferimos en este trabajo es la de DDPM/DDIM, cuya
formulación es la del \emph{denoising} probabilístico y no la de la
dinámica de Langevin sobre scores.

\todo{Si el tribunal lo requiere, desarrollar formalmente la
equivalencia entre la predicción de ruido $\epsilon_\theta$ y el score
$\nabla_{x_t}\log q(x_t)$, mostrando el factor de reescalado
$-1/\sqrt{1-\bar\alpha_t}$. Decidir el nivel de profundidad
matemática deseado para la defensa.}

\todo{Revisar / confirmar la atribución histórica: Sohl-Dickstein
(2015) como origen termodinámico, Song \& Ermon (2019) como score-
matching, Song et al.\ (2021, SDE) como unificación, Ho et al.\ (2020,
DDPM) como popularización. Añadir cualquier referencia adicional que
el tutor considere canónica.}

% ---------------------------------------------------------------------
\section{Principales aplicaciones}
\label{sec:aplicaciones-difusion}

El interés teórico se vio rápidamente respaldado por resultados
empíricos que situaron a la difusión por delante de los GAN en
síntesis de imágenes~\cite{dhariwal2021beatgans}. Desde entonces, el
abanico de aplicaciones se ha ampliado de forma notable. A modo de
panorámica:

\begin{itemize}
    \item \textbf{Síntesis de imágenes incondicional y condicional},
    incluyendo la generación texto-a-imagen a gran escala, viable
    gracias a los modelos de difusión latente que operan en un espacio
    comprimido en lugar de en el espacio de píxeles~\cite{rombach2022ldm}.
    \item \textbf{Problemas inversos en imagen}: super-resolución,
    \textit{inpainting}, \textit{deblurring} y \emph{denoising}, donde
    se busca recuperar una señal limpia a partir de una observación
    degradada.
    \item \textbf{Imagen médica}, que es el contexto de este trabajo:
    reconstrucción, reducción de ruido y síntesis de modalidades, con
    aplicaciones específicas a PET~\cite{scorepet2024,ddpmpet2023} que
    motivan directamente la contribución de esta memoria.
    \item \textbf{Otros dominios}: generación de audio y voz, de vídeo,
    de estructuras moleculares y de datos científicos, donde la misma
    receta de \emph{corromper y aprender a revertir} se traslada con
    pocas modificaciones.
\end{itemize}

La aplicación que vertebra este trabajo —la reconstrucción de PET de
baja dosis— se sitúa en la intersección entre los problemas inversos
en imagen y la imagen médica, y se desarrolla en detalle en el
Capítulo~\ref{cap:pet}.

\todo{Ajustar las referencias de aplicaciones a las que finalmente se
citen en la introducción para evitar duplicidades, y añadir las que
falten (p.\,ej.\ audio/vídeo si se quiere mencionar explícitamente).
Verificar la etiqueta \texttt{cap:pet} del capítulo de PET.}

% ---------------------------------------------------------------------
\section{Formulación del proceso de difusión}
\label{sec:formulacion-difusion}

En el contexto de generación de datos, se parte de una distribución
de probabilidad $p(x)$ que representa los datos de entrenamiento.
En el caso de las imágenes, podemos pensar en esta distribución
como una representación de todas las imágenes naturales; por
simplicidad, vamos a pensar en un subconjunto de imágenes
conteniendo únicamente caras de personas. Esta distribución es
tan compleja que no podemos encontrar una sola expresión que la
describa por completo. Igualmente, sin una fórmula específica,
buscamos generar imágenes nuevas, lo cual es equivalente a
muestrear puntos de esta distribución subyacente. El reto es
encontrar una forma de crear nuevos ejemplos sin tener una manera
comprensiva de encontrar la distribución. Los modelos de difusión
resuelven este problema con un enfoque completamente diferente.

Se construye un proceso de difusión que transforma cualquier
distribución en una distribución normal $\mathcal{N}(0,1)$. La
distribución condicional $q$ se define para ir de un paso al
siguiente; la distribución condicional $p$ se define para ir hacia
atrás, en la dirección opuesta. En el proceso hacia delante todos
los parámetros están fijos, pero en el proceso hacia atrás el
objetivo es encontrar los mejores parámetros $\theta$ que ayuden
a quitar el ruido de manera efectiva. Se entrenará una red neuronal
precisamente con el fin de encontrar estos parámetros.

¿Cómo se entrena esta red neuronal? Para entrenarla minimizaremos
la log-verosimilitud (negativa) de la probabilidad de producir una
muestra con nuestro modelo. De manera más sencilla, esto significa
encontrar el conjunto de parámetros $\theta$ que maximizan la
probabilidad de generar muestras reales $x_0$ de la distribución
de nuestros datos usando la red neuronal.

La probabilidad conjunta que describe la distribución de todas
las variables desde $x_1$ hasta $x_T$ dado $x_0$ representa el
proceso hacia delante por completo:
\begin{equation}
    q(x_1, \ldots, x_T \mid x_0)
    = q(x_1 \mid x_0)\, q(x_2 \mid x_1, x_0) \, \ldots\,
      q(x_T \mid x_{T-1}, \ldots, x_0).
\end{equation}

El proceso hacia delante (\textit{forward process}) del modelo de
difusión es una cadena de Markov, lo cual significa que cada paso
del proceso de adición de ruido depende exclusivamente del paso
anterior, simplificando las probabilidades condicionales:
\begin{equation}
    q(x_1, \ldots, x_T \mid x_0)
    = q(x_1 \mid x_0)\, q(x_2 \mid x_1) \, \ldots\,
      q(x_T \mid x_{T-1}).
\end{equation}

Podemos establecer los procesos hacia delante y hacia atrás como:
\begin{align}
    q(x_{1:T} \mid x_0)
        &= \prod_{t=1}^{T} q(x_t \mid x_{t-1}), \\
    p_\theta(x_{0:T})
        &= p_\theta(x_T) \prod_{t=1}^{T} p_\theta(x_{t-1} \mid x_t).
\end{align}

\paragraph{Referencias fundacionales.}
El paper que popularizó los modelos de difusión, \textit{Denoising
Diffusion Probabilistic Models}~\cite{ho2020ddpm}, sentó la base
formal y empírica sobre la que se asienta el resto del trabajo.
La referencia más temprana conocida a los modelos de difusión,
\textit{Deep Unsupervised Learning using Nonequilibrium
Thermodynamics}~\cite{sohldickstein2015}, plantea el marco
termodinámico que dio origen al enfoque.


% =====================================================================
% Capítulo 3 — Experimento preliminar: mariposas
% =====================================================================
\chapter{Generando mariposas: modelos de difusión en acción}
\label{cap:mariposas}

Como primer experimento práctico, se entrena un modelo DDPM
(\textit{Denoising Diffusion Probabilistic Model}) sobre un
conjunto de imágenes de mariposas, con el objetivo de generar
muestras sintéticas nuevas. Este ejercicio sirve como banco de
pruebas para comprender el proceso de entrenamiento e inferencia
antes de abordar el problema de reconstrucción de imágenes PET.

\section{Datos y preprocesamiento}
\label{sec:mariposas-datos}

Se utiliza el subconjunto \texttt{huggan/smithsonian\_butterflies\_subset}
del Instituto Smithsoniano. Cada imagen se redimensiona a
$256 \times 256$ píxeles preservando su relación de aspecto: el
lado más largo se reescala al tamaño objetivo, y el lado más corto
se rellena con relleno blanco simétrico hasta alcanzar una imagen
cuadrada (operación \texttt{ResizeToFit} + \texttt{PadToSquare} en
el código fuente). Sobre el resultado se aplica volteo horizontal
aleatorio como aumento de datos, y se normaliza al rango $[-1, 1]$
mediante la transformación
\begin{equation}
    x' = \frac{x - 0.5}{0.5}.
\end{equation}
El conjunto se carga en lotes de orden aleatorio en cada época.

Una iteración previa de este experimento siguió literalmente el
\textit{tutorial} oficial de la librería \texttt{diffusers}, con
imágenes reescaladas directamente a $128 \times 128$ píxeles. Ese
reescalado directo (sin preservar relación de aspecto ni rellenar
para mantener la geometría original) introduce una deformación
sistemática del contenido, visible en las muestras generadas
(\Cref{fig:mariposas-128}). El paso a $256 \times 256$ con
\texttt{ResizeToFit} + \texttt{PadToSquare} corrige esta patología
y es el preprocesamiento empleado en el resto del capítulo.

\section{Proceso de difusión hacia delante}
\label{sec:mariposas-forward}

El proceso de difusión hacia delante $q$ corrompe progresivamente
una imagen limpia $\mathbf{x}_0$ añadiendo ruido gaussiano a lo
largo de $T = 1000$ pasos discretos. En cada paso $t$, la
distribución de la imagen ruidosa condicionada a la imagen
anterior es
\begin{equation}
    q(\mathbf{x}_t \mid \mathbf{x}_{t-1})
    = \mathcal{N}\!\left(\mathbf{x}_t;\, \sqrt{1 - \beta_t}\,
      \mathbf{x}_{t-1},\, \beta_t \mathbf{I}\right),
\end{equation}
donde $\{\beta_t\}_{t=1}^{T}$ es un calendario de varianza
predefinido. Gracias a la propiedad de reproducibilidad de la
distribución gaussiana, este proceso admite una forma cerrada
que permite muestrear $\mathbf{x}_t$ directamente desde
$\mathbf{x}_0$ en un único paso:
\begin{equation}
    q(\mathbf{x}_t \mid \mathbf{x}_0)
    = \mathcal{N}\!\left(\mathbf{x}_t;\, \sqrt{\bar{\alpha}_t}\,
      \mathbf{x}_0,\, (1 - \bar{\alpha}_t)\mathbf{I}\right),
\end{equation}
donde $\alpha_t = 1 - \beta_t$ y $\bar{\alpha}_t = \prod_{s=1}^{t}
\alpha_s$. En la práctica, esto equivale al muestreo por
reparametrización
\begin{equation}
    \mathbf{x}_t
    = \sqrt{\bar{\alpha}_t}\,\mathbf{x}_0
    + \sqrt{1 - \bar{\alpha}_t}\,\boldsymbol{\varepsilon},
    \qquad
    \boldsymbol{\varepsilon} \sim \mathcal{N}(\mathbf{0}, \mathbf{I}).
\end{equation}

\section{Arquitectura del modelo}
\label{sec:mariposas-arquitectura}

La red neuronal $\boldsymbol{\varepsilon}_\theta(\mathbf{x}_t, t)$,
encargada de predecir el ruido, es una UNet2D con bloques ResNet
y atención espacial. La arquitectura presenta seis niveles de
resolución con canales $(128, 128, 256, 256, 512, 512)$ y dos
capas ResNet por bloque. La atención multi-cabeza se incorpora
únicamente en el quinto nivel, tanto en el camino descendente
(\texttt{AttnDownBlock2D}) como en el ascendente
(\texttt{AttnUpBlock2D}), donde la resolución espacial es
suficientemente reducida para que el coste cuadrático de la
atención sea asumible.

El paso temporal $t$ se inyecta en cada bloque ResNet mediante
\textit{embeddings} sinusoidales, análogos a los empleados en el
Transformer original, de modo que el modelo aprende una política
de \textit{denoising} condicional en el nivel de ruido.

\section{Objetivo de entrenamiento}
\label{sec:mariposas-loss}

En lugar de predecir directamente $\mathbf{x}_0$, el modelo sigue
la formulación simplificada de Ho et al.~\cite{ho2020ddpm} y
aprende a predecir el ruido $\boldsymbol{\varepsilon}$ a partir
de la imagen ruidosa. La función de pérdida es el error cuadrático
medio entre el ruido real y el predicho:
\begin{equation}
    \mathcal{L}
    = \mathbb{E}_{t \sim \mathcal{U}[1,T],\, \mathbf{x}_0,\,
                   \boldsymbol{\varepsilon}}
      \!\left[\left\|\boldsymbol{\varepsilon}
        - \boldsymbol{\varepsilon}_\theta\!\left(\mathbf{x}_t, t\right)
        \right\|^2\right].
\end{equation}
En cada iteración se muestrea $t$ uniformemente en
$\{1, \ldots, 1000\}$, se construye $\mathbf{x}_t$ mediante la
ecuación de muestreo directo, y el gradiente se propaga únicamente
a través de la predicción del modelo.

\section{Configuración del entrenamiento}
\label{sec:mariposas-config}

El modelo se optimiza con AdamW con tasa de aprendizaje inicial
$\eta = 10^{-4}$ y un planificador cosenoidal con 500 pasos de
calentamiento (\textit{warmup}). El entrenamiento se lleva a cabo
durante 100 épocas con precisión mixta \texttt{fp16} para reducir
el consumo de memoria, y se recortan los gradientes a norma
máxima 1.0 para garantizar la estabilidad numérica. Toda esta
fase preliminar se ejecuta \textbf{localmente} sobre un MacBook
Pro M4 Pro con aceleración MPS, suficiente para entrenar el
modelo a $256 \times 256$ en un plazo razonable y validar el
flujo de entrenamiento e inferencia antes de pasar a la tarea de
reconstrucción PET (véase \Cref{sec:pet-infraestructura}).

\section{Proceso de generación (inferencia)}
\label{sec:mariposas-inferencia}

Para generar una nueva imagen, se parte de ruido puro
$\mathbf{x}_T \sim \mathcal{N}(\mathbf{0}, \mathbf{I})$ y se
desnuda iterativamente aplicando el paso inverso aprendido
durante $T$ iteraciones. En cada paso $t$, el \textit{scheduler}
DDPM recupera $\mathbf{x}_{t-1}$ como
\begin{equation}
    \mathbf{x}_{t-1}
    = \frac{1}{\sqrt{\alpha_t}}
      \!\left(\mathbf{x}_t
        - \frac{\beta_t}{\sqrt{1 - \bar{\alpha}_t}}\,
        \boldsymbol{\varepsilon}_\theta(\mathbf{x}_t, t)\right)
    + \sigma_t\,\mathbf{z},
    \qquad
    \mathbf{z} \sim \mathcal{N}(\mathbf{0}, \mathbf{I}).
\end{equation}
Tras los $T = 1000$ pasos de \textit{denoising}, $\mathbf{x}_0$
constituye la imagen sintética final. La implementación permite
fijar una semilla aleatoria para garantizar reproducibilidad o
variarla libremente para explorar el espacio generativo del modelo
entrenado.

\section{Aprendizajes y resultados}
\label{sec:mariposas-aprendizajes}

Los modelos de difusión son buenos generalizando, ya que es más
complicado que puedan memorizar los ejemplos de entrenamiento,
dado que van a ver las mismas imágenes pero con una cantidad de
ruido completamente diferente entre una época y otra. Se ven
forzados a aprender la distribución subyacente de las imágenes
que analizan, en este caso aprendiendo cómo es una mariposa.
Aprender a predecir la cantidad de ruido que hay en una imagen
es matemáticamente equivalente a predecir la imagen de una
mariposa a partir de una imagen ruidosa.

Se ha observado que la pérdida durante el entrenamiento decrece
muy rápidamente en las primeras épocas, y después se estabiliza
a lo largo de todo el proceso de entrenamiento. Esto no significa
que el modelo deje de aprender una vez se estabiliza la pérdida,
sino que aprende a predecir la cantidad de ruido que hay en las
imágenes más rápidamente; pero los parámetros se siguen
actualizando y aprendiendo con más certeza la distribución
subyacente de imágenes de mariposas.

Para ilustrar este fenómeno comparamos las muestras generadas por
el modelo en dos momentos opuestos del entrenamiento: las primeras
20 épocas y las últimas 20 épocas (cada panel corresponde al
\textit{checkpoint} guardado al final de la época indicada).

La \Cref{fig:mariposas-tempranas} recoge las muestras de las
primeras 20 épocas, concretamente las correspondientes a los
\textit{checkpoints} de las épocas~10 y~20. En esta fase temprana
las mariposas apenas son reconocibles: aparecen como manchas de
color difusas, sin contornos definidos ni simetría, y con
texturas dominadas por artefactos de alta frecuencia. El modelo ya
ha captado la paleta cromática general del conjunto de datos
(tonos cálidos sobre fondos vegetales), pero todavía no ha
aprendido la estructura espacial que caracteriza a una mariposa.

\begin{figure}[h!]
    \centering
    \includegraphics[width=0.35\linewidth]{img/generated_butterflies/256pix_0009.png}\hspace{1em}
    \includegraphics[width=0.35\linewidth]{img/generated_butterflies/256pix_0019.png}
    \caption{Muestras de las primeras 20 épocas a resolución
    $256 \times 256$: \textit{checkpoints} al final de las
    épocas~10 (izquierda) y~20 (derecha). Las muestras son borrosas
    y carecen de estructura coherente, aunque ya reproducen la
    paleta cromática del conjunto de entrenamiento.}
    \label{fig:mariposas-tempranas}
\end{figure}

En contraste, la \Cref{fig:mariposas-tardias} muestra las muestras
de las últimas 20 épocas, correspondientes a los
\textit{checkpoints} de las épocas~90 y~100. A pesar de que la
pérdida es prácticamente la misma que en épocas más tempranas, la
calidad perceptual ha mejorado de forma notable: la red genera
mariposas con siluetas reconocibles, simetría bilateral aproximada
y patrones de color coherentes con la distribución de
entrenamiento. Se observa que las imágenes empiezan a ser buenas
en torno a la época~90 y se estabilizan en la época~100.

\begin{figure}[h!]
    \centering
    \includegraphics[width=0.35\linewidth]{img/generated_butterflies/256pix_0089.png}\hspace{1em}
    \includegraphics[width=0.35\linewidth]{img/generated_butterflies/256pix_0099.png}
    \caption{Muestras de las últimas 20 épocas a resolución
    $256 \times 256$: \textit{checkpoints} al final de las
    épocas~90 (izquierda) y~100 (derecha). Las muestras exhiben
    siluetas y simetrías reconocibles y patrones de color
    consistentes con la distribución de entrenamiento.}
    \label{fig:mariposas-tardias}
\end{figure}

\paragraph{Comparación en inferencia: $128 \times 128$ frente a
$256 \times 256$.}
Una vez analizada la evolución del entrenamiento, comparamos la
calidad de las muestras obtenidas en inferencia con las dos
resoluciones consideradas. La \Cref{fig:mariposas-128} muestra dos
ejemplos generados por la iteración previa del experimento,
entrenada con reescalado directo a $128 \times 128$ sin preservar
la relación de aspecto. Aunque la red aprende ciertas
regularidades de color y forma, las mariposas aparecen
visiblemente deformadas, lo cual ilustra el impacto del
preprocesamiento sobre la calidad final.

\begin{figure}[h!]
    \centering
    \includegraphics[width=0.35\linewidth]{img/generated_butterflies/128pix_sample1.jpeg}\hspace{1em}
    \includegraphics[width=0.35\linewidth]{img/generated_butterflies/128pix_sample2.jpeg}
    \caption{Muestras en inferencia a resolución $128 \times 128$
    (reescalado directo sin \textit{padding}). Las muestras
    presentan deformaciones sistemáticas heredadas del
    preprocesamiento.}
    \label{fig:mariposas-128}
\end{figure}

Por su parte, la \Cref{fig:mariposas-256} recoge dos ejemplos
generados en inferencia a resolución $256 \times 256$. Frente a la
línea base anterior, estas muestras presentan siluetas mucho más
nítidas y proporciones realistas, sin las deformaciones
sistemáticas de la versión a $128 \times 128$. El aumento de
resolución, combinado con un preprocesamiento que preserva la
relación de aspecto, se traduce en mariposas con mayor nivel de
detalle y mejor definición de los patrones alares.

\begin{figure}[h!]
    \centering
    \includegraphics[width=0.35\linewidth]{img/generated_butterflies/256pix_sample1.jpeg}\hspace{1em}
    \includegraphics[width=0.35\linewidth]{img/generated_butterflies/256pix_sample2.jpeg}
    \caption{Muestras en inferencia a resolución $256 \times 256$.
    Frente a la línea base de $128 \times 128$, las muestras son
    más nítidas, con proporciones realistas y mayor nivel de
    detalle en los patrones alares.}
    \label{fig:mariposas-256}
\end{figure}


% =====================================================================
% Capítulo 4 — Reconstrucción de PET de dosis baja
% =====================================================================
\chapter{Modelos de difusión para reconstrucción de escaneos PET}
\label{cap:pet}

\section{Caracterización del ruido en PET}
\label{sec:pet-ruido}

Una idea clave para desarrollar este trabajo es utilizar los
escaneos de máxima calidad y entrenar un modelo de difusión de
tal forma que las funciones de adición de ruido imiten el tipo
de ruido presente en los escaneos de dosis reducida. Típicamente,
los datos PET están muestreados con ruido de Poisson de alta
varianza.

Un malentendido frecuente es que el modelo de ruido de Poisson
del PET es incompatible con el ruido gaussiano artificial empleado
en el proceso de difusión. Esto no es un problema. En general,
el proceso de difusión establece una correspondencia entre una
distribución útil de imágenes médicas (con cierto ruido inherente)
y una distribución conocida. Esta distribución conocida se elige
como gaussiana por sus propiedades matemáticas convenientes,
aunque otras opciones son posibles.

\section{Planteamiento de la solución}
\label{sec:pet-planteamiento}

A partir de los fundamentos repasados anteriormente y de las
propuestas recogidas en el estado del arte
(\Cref{sec:estado-del-arte}), se diseña la solución para la tarea
de reconstrucción de PET de dosis baja con modelos de difusión.
El diseño persigue dos objetivos: ser \textbf{defendible
científicamente} —cada decisión está justificada por la literatura
recogida o por los fundamentos teóricos repasados— y
\textbf{ejecutable sobre el hardware disponible}. El
entrenamiento se ejecutó sobre una NVIDIA RTX PRO 4500 en la
plataforma cloud RunPod, tras una fase preliminar de validación
del flujo sobre el equipo local empleado en el experimento de
mariposas —un MacBook Pro M4 Pro con 24~GB de memoria unificada y
aceleración MPS, sin acceso a CUDA— (véase
\Cref{sec:pet-infraestructura}).

La reconstrucción de PET de dosis baja se formula como un
problema inverso: dada una observación $\mathbf{y}$ (escaneo de
dosis reducida), se busca recuperar $\mathbf{x}$ (escaneo
equivalente a dosis completa). En la literatura coexisten dos
familias de soluciones basadas en difusión, ambas representadas
en el estado del arte:

\begin{enumerate}[leftmargin=*]
    \item \textbf{Difusión condicional supervisada.} Se entrena
    directamente la distribución posterior $p(\mathbf{x} \mid
    \mathbf{y})$, incorporando $\mathbf{y}$ como entrada de
    condicionamiento de la red de \textit{denoising}. Es la
    formulación más simple y aprovecha íntegramente los 371
    pares registrados disponibles.

    \item \textbf{Prior no condicional con consistencia de
    medición.} Se entrena un modelo $p(\mathbf{x})$ sobre los
    escaneos de dosis completa exclusivamente; en inferencia,
    la trayectoria inversa del muestreo se guía hacia
    consistencia con la observación $\mathbf{y}$ mediante un
    término de verosimilitud~\cite{chung2023dps}. Esta es la
    línea identificada como estado del arte para el PET de
    dosis baja por la revisión de junio de 2024 del
    \textit{British Journal of Radiology, Artificial
    Intelligence}~\cite{bjrai2024}.
\end{enumerate}

Se decide construir \textbf{ambos enfoques en paralelo} y
compararlos sobre el mismo conjunto de test. La comparación
entre el supervisado y el no condicional con guiado constituye
uno de los resultados centrales que la memoria persigue reportar.

\section{Datos: caracterización y preprocesamiento}
\label{sec:pet-datos}

El conjunto disponible consta de \textbf{371 pares} de volúmenes
en formato NIfTI: para cada paciente se dispone de la versión
adquirida con dosis completa y la versión con dosis reducida a un
veinteavo. Cada volumen tiene forma $(440, 440, 644)$ con vóxeles
isotrópicos de $1{,}65$~mm. Las matrices afines (\texttt{affine})
coinciden entre los volúmenes pareados, lo cual confirma que la
reducción de dosis se aplica sobre el mismo grid espacial sin
re-muestreo intermedio.

Las estadísticas relevantes para el preprocesamiento, calculadas
sobre el volumen de referencia \texttt{01122021\_1}, son:

\begin{itemize}[leftmargin=*]
    \item La distribución de intensidades en \textit{foreground}
    es muy asimétrica: mediana de 251 cuentas, percentil 99 en
    torno a 9.600, percentil 99,9 en torno a 19.000 y máximo
    cercano a $1{,}3 \times 10^{5}$.
    \item Aproximadamente el 18\,\% de los vóxeles del volumen
    son \textit{foreground} (valor estrictamente positivo).
    \item En la región de \textit{foreground}, la transformación
    $\log(1+x)$ aproxima una distribución $\mathcal{N}(4{,}7,\,
    2{,}6)$, lo cual sugiere que una compresión logarítmica
    produce un rango bien centrado para el modelo.
    \item La desviación típica del residuo
    \texttt{low\,-\,full} restringido al \textit{foreground} es
    del orden del 31\,\% de la media de señal, lo cual
    cuantifica el nivel de degradación introducido por la
    reducción de dosis.
\end{itemize}

Las decisiones de preprocesamiento son las siguientes.

\paragraph{Unidad espacial.}
Se trabaja con \textit{slices} axiales en 2D puro. Un U-Net 3D
sobre volúmenes de $440^3$ es inviable con 24~GB de memoria
unificada; las variantes 2.5D (canales adicionales con
\textit{slices} contiguos) se reservan como extensión natural si
la formulación 2D muestra inconsistencias inter-\textit{slice}.

\paragraph{Recorte y reescalado.}
Para cada volumen se calcula una \textit{bounding box} única a
partir del \textit{foreground} de la versión de \textbf{dosis
baja}, y se aplica idéntica al volumen pareado de dosis completa,
garantizando que ambos comparten exactamente el mismo recorte.
Derivar el recorte de la dosis baja —y no de la completa— evita
cualquier fuga de información del objetivo hacia la entrada: la
\textit{bounding box}, la escala de normalización $M$ y el filtrado
de \textit{slices} se calculan todos a partir de la observación de
dosis baja, la única señal disponible en inferencia para un
volumen no visto. Cada
\textit{slice} axial recortado se redimensiona a $256 \times 256$
píxeles mediante interpolación bilineal. Esta resolución
concentra la capacidad del modelo en la anatomía relevante y
reduce la carga computacional respecto al $440 \times 440$ nativo.

La resolución de trabajo $256^2$ es menor que la nativa
($440 \times 440$). No es posible entrenar directamente sobre
$440^2$: la arquitectura U-Net empleada tiene seis bloques —cinco
operaciones de \textit{down}/\textit{upsampling}, $2^5 = 32$—, de
modo que el lado de la imagen de entrada debe ser divisible entre
32, y $440$ no lo es. Tamaños como $384$ o $416$ —más cercanos a
la calidad nativa y divisibles entre 32— sí serían admisibles,
pero no se han explorado: el coste de cualquier fase del
entrenamiento crece de forma aproximadamente cuadrática con el
lado de la imagen, y ya resulta elevado a $256^2$. La elección de
$256^2$ equilibra fidelidad anatómica y coste computacional dentro
del alcance del proyecto.

\paragraph{Normalización per-volumen estabilizadora de varianza.}
Las cuentas se mapean al rango $[-1, 1]$ mediante la
transformación
\begin{equation}
    \mathbf{x}' = 2 \cdot
        \frac{\operatorname{arcsinh}(\mathbf{x}/k)}
             {\operatorname{arcsinh}(M/k)}
        - 1,
\end{equation}
donde $k = 10$ actúa como codo suave entre el régimen lineal y
el logarítmico, y $M$ es el percentil 99,5 del \textbf{volumen de
dosis baja}. El mismo $M$ derivado de la dosis baja se aplica
también al volumen pareado de dosis completa, de modo que ambos
comparten escala y siguen siendo directamente comparables tras la
normalización. Calcular $M$ exclusivamente sobre la dosis baja es
esencial para no introducir fuga del objetivo: $M$ debe poder
recomputarse en inferencia a partir de la única señal disponible,
el escaneo de dosis baja de un volumen no visto. La función
$\operatorname{arcsinh}$ es la transformación canónica
estabilizadora de varianza para señales positivas de cola pesada
con estructura aproximadamente Poisson: se comporta linealmente
cerca de cero y logarítmicamente para valores grandes, sin la
patología de $\log(0)$. Los parámetros $(M, k)$ se almacenan por
volumen para invertir la transformación durante la evaluación.

\paragraph{Elección del percentil de referencia.}
La elección de $M$ como percentil 99,5 —y no como el máximo
(percentil 100)— es deliberada. Si se fijara $M$ al valor más alto
del volumen, un único vóxel atípico muy brillante podría
desviarse mucho de la distribución y aplanar el resto de la
imagen, perdiendo detalle en la masa de tejido. Por otra parte,
los modelos de difusión asumen que los datos viven en $[-1, +1]$:
es para esa banda para la que está ajustado el calendario de
ruido. Los vóxeles que quedan bastante por encima de $+1$ están
efectivamente fuera de distribución para el proceso de ruido hacia
delante, y el modelo tiende a reconstruirlos peor (saturación
suave / subestimación de la intensidad máxima). No se
\textit{clippean} esos valores —dejarlos sobresalir por encima de
$+1$ mantiene exacta la inversión de la transformación—, pero su
proporción debe ser pequeña. Un análisis del porcentaje de
\textit{foreground} mapeado por encima de $+1$ para distintos
percentiles arroja:
\begin{center}
\begin{tabular}{lc}
\toprule
Percentil & \% \textit{foreground} mapeado por encima de $+1$ \\
\midrule
90{,}0 & 10{,}00\,\% \\
95{,}0 & 5{,}00\,\% \\
99{,}0 & 1{,}00\,\% \\
\textbf{99{,}5} & \textbf{0{,}50\,\%} (por defecto) \\
99{,}9 & 0{,}10\,\% \\
\bottomrule
\end{tabular}
\end{center}
Bajar el percentil dedicaría más del rango $[-1, +1]$ a la masa
de vóxeles —lo cual podría mejorar las métricas estructurales en
tejido blando (SSIM, contraste de intensidad media)— a costa de
empeorar la fidelidad de intensidad máxima en regiones calientes
y las métricas de preservación de intensidad en espacio de
cuentas. El percentil 99,5 deja fuera únicamente un 0,5\,\% de
\textit{foreground} —\textit{outliers} muy brillantes—, lo cual
preserva la distribución original de los datos a la vez que los
normaliza de forma efectiva.

\paragraph{Filtrado de \textit{slices}.}
Tras el recorte y la normalización, se conservan únicamente
aquellos \textit{slices} axiales cuyo \textit{foreground} supera
un umbral mínimo de ocupación (típicamente $1\%$ de vóxeles por
encima de un valor bajo). Los extremos vacíos del FOV se descartan
para no penalizar al modelo con cientos de \textit{slices}
triviales por volumen.

\paragraph{Ausencia de aumento de datos.}
El PET presenta asimetría clínica relevante —corazón en el lado
izquierdo, hígado en el derecho, etcétera—; un volteo horizontal
aleatorio induciría una lateralidad incorrecta en el modelo
aprendido. El tamaño efectivo del conjunto de entrenamiento
(aproximadamente $1{,}5 \times 10^{5}$ \textit{slices} con
\textit{foreground}) hace innecesario recurrir a la
augmentación para evitar sobreajuste.

\paragraph{Particiones.}
Los 371 volúmenes se dividen al nivel de fichero en proporción
$80 / 10 / 10$ ($\sim 296 / 37 / 38$ volúmenes) con semilla fija,
y la asignación se persiste en un fichero JSON compartido por
ambos \textit{pipelines}. Se asume que cada fichero corresponde
a un paciente independiente; el patrón de los \textit{timestamps}
en los nombres es consistente con un proceso de reconstrucción
por lotes en el operador, no con \textit{frames} dinámicos del
mismo paciente.

\paragraph{Caché en disco.}
El preprocesamiento se ejecuta una única vez en modo
\textit{offline}; cada \textit{slice} procesado se almacena como
tensor PyTorch en formato \texttt{.pt} con precisión
\texttt{float16} para evitar la decodificación repetida de
NIfTI en cada época. El tamaño total estimado del caché es de
aproximadamente 60~GB.

\section{Arquitectura y paradigma de entrenamiento}
\label{sec:pet-arquitectura}

La arquitectura del U-Net es idéntica a la empleada en el modelo
de mariposas, modificando exclusivamente los canales de entrada
y salida. Las dimensiones y la topología se conservan:

\begin{itemize}[leftmargin=*]
    \item \texttt{block\_out\_channels = (128, 128, 256, 256, 512,
    512)}, lo cual produce un modelo de aproximadamente 114
    millones de parámetros (idéntico recuento al modelo de
    mariposas, dado que el primer canal convolucional contribuye
    una fracción despreciable al total).
    \item \texttt{layers\_per\_block = 2}.
    \item Bloque de atención multi-cabeza en el quinto nivel del
    camino descendente y ascendente (resolución espacial
    $16 \times 16$ tras cinco \textit{downsamplings} del
    \textit{input} de $256 \times 256$).
\end{itemize}

Los canales de entrada y salida dependen del \textit{pipeline}:

\begin{itemize}[leftmargin=*]
    \item \textbf{Pipeline A (supervisado):} \texttt{in\_channels = 2}
    —canal de difusión $\mathbf{x}_t$ y canal de condicionamiento
    $\mathbf{y}$—, \texttt{out\_channels = 1}.
    \item \textbf{Pipeline B (prior no condicional):}
    \texttt{in\_channels = 1}, \texttt{out\_channels = 1}.
\end{itemize}

El paradigma de difusión combina tres elecciones que mejoran
sobre la línea base ($\varepsilon$-prediction + \textit{schedule}
lineal) del modelo de mariposas:

\paragraph{Predicción de velocidad (\textit{v-prediction}).}
Salimans y Ho~\cite{salimans2022vprediction} proponen
reparametrizar la red para que prediga
$\mathbf{v}_t = \alpha_t\,\boldsymbol{\varepsilon} -
\sigma_t\,\mathbf{x}_0$ en lugar de $\boldsymbol{\varepsilon}$.
Esta parametrización tiene un comportamiento numérico estable
en los dos límites $t \to 0$ y $t \to T$, lo cual evita
inestabilidades de entrenamiento en los \textit{timesteps} de
muy bajo o muy alto ruido. Su activación en \texttt{diffusers}
consiste únicamente en fijar
\texttt{prediction\_type="v\_prediction"} en el \textit{scheduler}.

\paragraph{\textit{Schedule} cosenoidal.}
Nichol y Dhariwal~\cite{nichol2021cosine} demuestran
empíricamente que un \textit{schedule} de varianza $\beta_t$ con
perfil cosenoidal produce una distribución de SNR a lo largo de
los \textit{timesteps} mejor adaptada a resoluciones medias que
el \textit{schedule} lineal original de Ho et al.~\cite{ho2020ddpm}.
Se mantienen $T = 1000$ \textit{timesteps} de entrenamiento.

\paragraph{Pérdida MSE uniforme sobre $\mathbf{v}$.}
La función de pérdida es
\begin{equation}
    \mathcal{L}
    = \mathbb{E}_{t,\, \mathbf{x}_0,\, \boldsymbol{\varepsilon}}
      \!\left[\left\|\mathbf{v}_t
        - \mathbf{v}_\theta(\mathbf{x}_t,\, \mathbf{c},\, t)
      \right\|^2\right],
\end{equation}
donde $\mathbf{c}$ representa el canal de condicionamiento
$\mathbf{y}$ en Pipeline~A y se omite en Pipeline~B. Se opta por
una ponderación uniforme sobre los \textit{timesteps}, dejando
el reponderado por SNR mínimo (\textit{Min-SNR-$\gamma$}, Hang et
al.~\cite{hang2023minsnr}) como mejora candidata si la
convergencia resulta insuficiente.

\paragraph{Hiperparámetros de optimización.}
Los hiperparámetros se ajustan a la realidad del hardware
finalmente empleado —la GPU NVIDIA RTX PRO 4500 sobre RunPod
detallada en \Cref{sec:pet-infraestructura}—:

\begin{itemize}[leftmargin=*]
    \item Optimizador \texttt{AdamW} con tasa de aprendizaje
    $\eta = 1{,}4 \times 10^{-4}$ y \textit{cosine warmup} de
    500 pasos. El valor de $\eta$ está escalado por
    $\sim\sqrt{32/16}$ sobre la línea base $10^{-4}$ empleada en
    el modelo de mariposas, para acompañar el incremento de
    \textit{batch} sin alterar la dinámica del entrenamiento.
    \item \textit{Exponential moving average} (EMA) sobre los
    pesos del U-Net con factor de decaimiento $0{,}9999$. Los
    pesos EMA son los que se evalúan y se guardan como
    \textit{checkpoint} definitivo; aportan típicamente $0{,}5$--$1$
    dB adicionales de PSNR sin coste de entrenamiento.
    \item Precisión mixta \texttt{bf16} sobre los \textit{tensor
    cores} Blackwell, prácticamente el doble de \textit{throughput}
    que \texttt{float32} sin riesgo apreciable de NaN. (En el
    \textit{fallback} local sobre MPS se emplea \texttt{float32}
    por las inestabilidades conocidas del \textit{backend} con
    \texttt{float16}.)
    \item \textit{Batch} de entrenamiento de 32 \textit{slices}
    con acumulación de gradiente unitaria, lo cual produce un
    \textit{batch} efectivo de 32. Sobre MPS se utiliza
    \textit{micro-batch} de 4 con acumulación de 4 pasos para
    obtener un \textit{batch} efectivo de 16, idéntico al del
    modelo de mariposas.
    \item Se entrena durante 30 épocas, lo cual equivale a
    aproximadamente $1{,}1 \times 10^{6}$ pasos de gradiente
    sobre las $\sim 1{,}5 \times 10^{5}$ \textit{slices} con
    \textit{foreground}.
\end{itemize}

\section{Pipeline~A — supervisado condicional}
\label{sec:pet-pipeline-a}

El condicionamiento se realiza mediante \textbf{concatenación por
canales}, el mecanismo dominante en la literatura supervisada de
PET de dosis baja con DDPMs. La entrada del U-Net es el tensor
de forma $2 \times 256 \times 256$ obtenido concatenando
$\mathbf{x}_t$ —el \textit{slice} de dosis completa con ruido
inyectado en el \textit{timestep}~$t$— y $\mathbf{y}$ —el
\textit{slice} pareado de dosis baja, normalizado a la misma
escala que $\mathbf{x}_0$—. La salida es la predicción
$\hat{\mathbf{v}}_t$ de tamaño $1 \times 256 \times 256$.

En inferencia, para cada \textit{slice} axial del volumen de test
se parte de $\mathbf{x}_T \sim \mathcal{N}(\mathbf{0}, \mathbf{I})$
concatenado con el correspondiente $\mathbf{y}$, y se aplica el
muestreador \textbf{DDIM}~\cite{song2021ddim} con 50 pasos de
\textit{denoising} y $\eta = 0$, lo cual hace el proceso
determinista para una semilla dada. Tras procesar todos los
\textit{slices} del volumen, éstos se ensamblan en un único
\textit{array} tridimensional, se aplica la transformación
inversa de la normalización per-volumen, y el resultado se
escribe en formato NIfTI con la misma matriz afín que el volumen
de entrada.

\section{Pipeline~B — prior no condicional con guiado de medición
(DPS)}
\label{sec:pet-pipeline-b}

El entrenamiento del Pipeline~B es idéntico al del Pipeline~A
pero \textbf{sin el canal de condicionamiento}: el U-Net aprende
$p(\mathbf{x})$ sobre los \textit{slices} de dosis completa
exclusivamente. La observación $\mathbf{y}$ no interviene en el
entrenamiento.

En inferencia se introduce la consistencia con la medición
mediante \textbf{Diffusion Posterior Sampling}~\cite{chung2023dps}
con operador directo identidad. En cada paso DDIM se ejecuta el
siguiente procedimiento:

\begin{enumerate}[leftmargin=*]
    \item La red predice $\hat{\mathbf{v}}_t$ a partir del estado
    actual $\mathbf{x}_t$.
    \item Se reconstruye la estimación del estado limpio
    mediante la inversa de la parametrización $\mathbf{v}$:
    \begin{equation}
        \hat{\mathbf{x}}_0
        = \alpha_t\,\mathbf{x}_t
        - \sigma_t\,\hat{\mathbf{v}}_t.
    \end{equation}
    \item Se aplica un paso de gradiente sobre $\mathbf{x}_t$ que
    penaliza el error de medición:
    \begin{equation}
        \mathbf{x}_t \leftarrow
            \mathbf{x}_t - \omega \cdot
            \nabla_{\mathbf{x}_t}
            \!\left[\frac{1}{2}
                \left\|\hat{\mathbf{x}}_0 - \mathbf{y}\right\|_2^2
            \right],
    \end{equation}
    donde $\omega$ es el escalar de guiado, único hiperparámetro
    \textit{tunable} del procedimiento.
    \item Se ejecuta el paso DDIM estándar para obtener
    $\mathbf{x}_{t-1}$ a partir del $\mathbf{x}_t$ ya corregido.
\end{enumerate}

La elección del operador directo identidad ($A = I$) y de una
varianza $\sigma$ uniforme es la aproximación más simple y
defendible en este contexto, ya que el PET de dosis baja carece
de una formulación limpia $\mathbf{y} = A\mathbf{x}$ —la
reducción de dosis combina re-muestreo de Poisson, reducción
del tiempo de escaneo y la posterior etapa de reconstrucción
algorítmica, ninguna de las cuales se modela explícitamente—.
La pérdida cuadrática en el espacio normalizado por
$\operatorname{arcsinh}$ equivale a una verosimilitud gaussiana
isotrópica sobre la representación estabilizada de varianza. La
sustitución de $\sigma$ uniforme por $\sigma(\mathbf{y})$ derivada
de la estadística de Poisson es una de las extensiones
identificadas como líneas de desarrollo futuro
(\Cref{sec:lineas-futuras}).

\section{Evaluación}
\label{sec:pet-evaluacion}

Sobre el conjunto de test ($\sim 38$ volúmenes) se reportan tres
familias de métricas.

\paragraph{Métricas de calidad de imagen.}
PSNR, SSIM y NRMSE se calculan por \textit{slice} y se agregan
por volumen mediante la media a lo largo de los \textit{slices}.
Se reportan dos variantes: (i) sobre el \textit{slice} completo
y (ii) restringidas a la máscara de \textit{foreground} del
\textit{slice}. La variante restringida es la clínicamente
relevante porque ignora el fondo vacío; la variante completa es
la convención estándar y permite la comparación con la literatura.

Estas tres métricas se calculan sobre el \textbf{espacio
normalizado} de las imágenes (la representación
$\operatorname{arcsinh}$ en $[-1, +1]$), y no sobre el espacio de
cuentas. La razón es que el PET tiene un rango dinámico altísimo:
el PSNR colapsa hacia el vóxel más brillante, por lo que evaluarlo
en cuentas equivaldría a medir casi exclusivamente cómo de bien se
reproduce el punto más caliente, siendo numéricamente insensible
a los errores del resto de la imagen; análogamente, los valores
atípicos inflarían el SSIM. Además, promediar a lo largo de
pacientes requiere una escala común para que las comparaciones
intra-volumen sean justas. La reformulación clave es que
$\operatorname{arcsinh}$ es monótona e invertible y se aplica por
igual a la reconstrucción y al \textit{ground truth}: no oculta el
error, sino que lo repondera para que cuenten los errores
relativos en todo el rango dinámico en lugar de los errores
absolutos en el pico. Así, PSNR/SSIM en espacio normalizado son
una comparación fiel del resultado final, bajo una ponderación de
intensidad en la que la métrica no está degenerada. La
preservación de intensidad cuantitativa, en cambio, sí se evalúa
en espacio de cuentas (véase a continuación).

\paragraph{Preservación de intensidad.}
Como métrica \textit{proxy} de la preservación de SUV —no se
dispone de metadatos de dosis inyectada ni peso del paciente para
computar SUV reales—, se reporta el error relativo entre la
\textbf{media} y el \textbf{máximo} de intensidad de
\textit{foreground} por volumen, calculado en \textbf{espacio
de cuentas original} (tras invertir la transformación
$\operatorname{arcsinh}$).

\paragraph{Evaluación cualitativa.}
Para $N = 5$ volúmenes representativos del conjunto de test se
generan figuras de cuatro paneles (dosis baja, dosis completa,
reconstrucción, residuo absoluto) que permiten un análisis visual
del comportamiento del modelo.

La comparación entre Pipeline~A y Pipeline~B se realiza sobre
exactamente los mismos volúmenes y los mismos \textit{slices},
lo cual garantiza que las métricas reportadas son directamente
comparables.

\subsection{Resultados}
\label{sec:pet-resultados}

La \Cref{tab:pet-supervisado} recoge las métricas del Pipeline~A
(supervisado condicional) sobre el conjunto de test, evaluado en
los \textit{checkpoints} de las épocas 50 y 100. La progresión
entre ambos \textit{checkpoints} es consistente en las métricas de
calidad de imagen: el PSNR sobre el \textit{slice} completo sube
de $34{,}52$ a $35{,}32$~dB, el SSIM mejora tanto en
\textit{foreground} como en el \textit{slice} completo, y el NRMSE
disminuye en ambas variantes. La preservación de intensidad media
en espacio de cuentas se mantiene en torno a un $-4$ a $-5$\,\% de
error relativo, lo cual indica una ligera subestimación
sistemática de la señal —coherente con la saturación suave
esperable en las regiones más calientes—.

\begin{table}[h!]
    \centering
    \begin{tabular}{lcc}
        \toprule
        Métrica & Época 50 & Época 100 \\
        \midrule
        PSNR \textit{foreground} (dB) & $29{,}94$ & $30{,}20$ \\
        PSNR \textit{slice} completo (dB) & $34{,}52$ & $35{,}32$ \\
        SSIM \textit{foreground} & $0{,}9638$ & $0{,}9652$ \\
        SSIM \textit{slice} completo & $0{,}9207$ & $0{,}9251$ \\
        NRMSE \textit{foreground} & $0{,}0375$ & $0{,}0364$ \\
        NRMSE \textit{slice} completo & $0{,}0213$ & $0{,}0194$ \\
        Error medio intensidad (\%) & $-4{,}18$ & $-4{,}75$ \\
        Error máximo intensidad (\%) & $-14{,}64$ & $-18{,}47$ \\
        \bottomrule
    \end{tabular}
    \caption{Métricas del Pipeline~A (supervisado condicional)
    sobre el conjunto de test, en los \textit{checkpoints} de las
    épocas 50 y 100. PSNR/SSIM/NRMSE se calculan en espacio
    normalizado; el error de intensidad, en espacio de cuentas.}
    \label{tab:pet-supervisado}
\end{table}

\todo{Añadir la columna del Pipeline~B (prior no condicional con
DPS) a la \Cref{tab:pet-supervisado} e insertar las figuras
representativas de cuatro paneles tras completar su evaluación.}

\section{Infraestructura de cómputo y migración a la nube}
\label{sec:pet-infraestructura}

\subsection{Plan inicial: cómputo local}

El planteamiento inicial, recogido en
\Cref{sec:pet-planteamiento}, contemplaba la ejecución completa
del entrenamiento y la inferencia en local sobre el mismo MacBook
Pro M4 Pro empleado en el experimento preliminar de mariposas:
24~GB de memoria unificada, PyTorch~2.10 con \textit{backend}
MPS, y un \textit{micro-batch} de 4 \textit{slices} con
acumulación de gradiente de 4 pasos en precisión \texttt{float32}
—la implementación de MPS presenta inestabilidades ocasionales
con \texttt{float16}—. Bajo estas condiciones, el coste estimado
era de aproximadamente \textbf{75~horas por pipeline} para
30~épocas y de \textbf{150~horas en total} para entrenar los dos
\textit{pipelines}, además de cerca de 2~horas por evaluación
completa del conjunto de test (38 volúmenes $\times$ $\sim 500$
\textit{slices} $\times$ 50 pasos DDIM).

\subsection{Motivación de la migración}

Aunque el cómputo local fue suficiente para el experimento
preliminar de mariposas, dos factores motivaron migrar la fase
PET a infraestructura en la nube. En primer lugar, el presupuesto
de cómputo era considerable: 150~horas de entrenamiento sostenido
acoplan el ordenador personal al proceso durante una semana
completa, con todas las interrupciones (cierre de tapa, suspensión
por inactividad, reinicios de sistema) traducidas a riesgo
operativo. En segundo lugar, el techo de rendimiento de MPS y el
uso forzado de \texttt{float32} fijaban una cota inferior al
tiempo total que la elección de mejor hardware podía relajar de
forma significativa.

\subsection{Entorno cloud: RunPod + RTX PRO 4500}

Se traslada la fase de entrenamiento a la plataforma
\textbf{RunPod}, un proveedor de GPU bajo demanda que permite
instanciar máquinas con GPU dedicada por horas y con tarifa
proporcional al uso. La GPU seleccionada es una \textbf{NVIDIA
RTX PRO 4500 (arquitectura Blackwell, 32~GB de VRAM, CUDA)}, lo
cual habilita varias mejoras frente al setup MPS:

\begin{itemize}[leftmargin=*]
    \item Reactivación de la precisión mixta (\texttt{bf16}
    sobre tensores Blackwell), prácticamente el doble de
    \textit{throughput} que \texttt{float32} sin riesgo
    apreciable de inestabilidad numérica.
    \item Aumento del \textit{train batch size} de 4 a 32 con
    acumulación de gradiente unitaria, ajustando la tasa de
    aprendizaje a $\eta = 1{,}4 \times 10^{-4}$
    ($\sim\sqrt{32/16}$ sobre el valor original) para mantener
    la dinámica de entrenamiento.
    \item Desacoplamiento total del proceso de entrenamiento
    respecto del ordenador personal: una vez lanzado, el
    \textit{pipeline} progresa de forma autónoma con
    independencia de la conectividad o estado del portátil.
\end{itemize}

\subsection{Flujo operativo}

La operativa de uso de la GPU remota descansa sobre cuatro
herramientas estándar de trabajo con sistemas Linux remotos:

\begin{description}[leftmargin=*]
    \item[Volumen de almacenamiento persistente.] RunPod permite
    asociar a la GPU un volumen de disco que sobrevive al
    apagado de la instancia. Sobre él se materializa el caché
    de \textit{slices} preprocesados ($\sim 60$~GB), los
    \textit{checkpoints} y los \textit{logs} de entrenamiento.
    De esta manera, detener la GPU para liberar costes cuando
    no se está usando no implica perder el estado del trabajo
    ni reprocesar los volúmenes NIfTI.
    \item[SSH.] La conexión al \textit{pod} se realiza por SSH
    sobre la dirección y puerto que RunPod expone tras lanzar la
    instancia. El cifrado del canal garantiza que las
    credenciales y los datos transferidos no sean accesibles
    desde la red intermedia.
    \item[\texttt{tmux}.] Cada sesión de entrenamiento se lanza
    dentro de una sesión \texttt{tmux} sobre el \textit{pod}
    remoto. \texttt{tmux} desacopla el ciclo de vida del proceso
    del ciclo de vida de la conexión SSH: cerrar la terminal
    local, perder la conexión a la red o reiniciar el portátil
    no interrumpen el entrenamiento, y reconectarse al
    \textit{pod} permite recuperar la sesión exactamente en el
    estado en el que se dejó.
    \item[Cursor / VSCode con Remote-SSH.] La edición y
    depuración del código se realiza con Cursor, un \textit{fork}
    de Visual Studio Code, utilizando la extensión \textit{Remote
    -- SSH}. El editor se ejecuta localmente pero opera sobre el
    sistema de archivos del \textit{pod} remoto a través del
    mismo canal SSH, lo cual proporciona la experiencia de
    desarrollo habitual (autocompletado, navegación de símbolos,
    terminal integrada) sobre la máquina con GPU sin necesidad
    de sincronizar copias locales del código.
\end{description}

\subsection{Implicación sobre la configuración descrita en la memoria}

La configuración numérica del entrenamiento (\textit{batch} de 32,
\texttt{bf16}, $\eta = 1{,}4 \times 10^{-4}$ con \textit{cosine
warmup} de 500 pasos, \texttt{gradient\_accumulation\_steps = 1},
30 épocas) refleja el setup CUDA en RunPod. El presupuesto en
horas-GPU sobre la RTX PRO 4500 se actualiza respecto al cálculo
inicial sobre MPS: \todo{registrar tiempo real de entrenamiento
por \textit{pipeline} y precio total en RunPod tras la
ejecución}.

\section{Estructura del código}
\label{sec:pet-codigo}

El código fuente se organiza como un único paquete
\texttt{pet\_reconstruction/src/} con infraestructura compartida
y módulos específicos por \textit{pipeline}:

\begin{verbatim}
pet_reconstruction/src/
  config.py                       # dataclasses, sub-configs por pipeline
  volume_io.py                    # NIfTI I/O, bbox, arcsinh + inverso
  splits.py                       # split 80/10/10 a nivel paciente
  preprocess.py                   # offline: NIfTI -> cache .pt
  data.py                         # PairedSliceDataset + FullDoseOnlyDataset
  metrics.py                      # PSNR/SSIM/NRMSE + preservacion intensidad
  visualize.py                    # figuras 4-panel
  model_supervised.py             # UNet builder, in_channels=2
  train_supervised.py             # bucle de entrenamiento Pipeline A
  reconstruct_supervised.py       # DDIM-50, salida NIfTI
  model_unconditional.py          # UNet builder, in_channels=1
  train_unconditional.py          # bucle de entrenamiento Pipeline B
  reconstruct_unconditional.py    # DDIM-50 + guiado DPS, salida NIfTI
  evaluate.py                     # inferencia + metricas -> CSV/JSON/figuras
  main.py                         # dispatcher
\end{verbatim}

Un preset \texttt{--smoke} (50 pacientes, 5 épocas, resolución
$128 \times 128$) permite verificar el flujo completo de extremo
a extremo en aproximadamente una hora antes de comprometerse a
los entrenamientos largos.
